\textbf{Basis Vector}
\begin{lstlisting}
template<typename T = int, int LG = 31>
struct Basis {
  T basis[LG]; int sz;
  vector<int> who;
  Basis() {
    fill(basis, basis + LG, 0), sz = 0;
    who.clear();
  }
  void reset() {
    fill(basis, basis + LG, 0), sz = 0;
    who.clear();
  }
  bool insert(T mask, int idx = -1) {
    for (int i = LG - 1; i >= 0; --i) {
      if (mask >> i & 1) {
        if (basis[i]) mask ^= basis[i];
        else {
          basis[i] = mask, ++sz;
          who.push_back(idx);
          return true;
        }
      }
    }
    return false;
  }
  T maxXor() {
    T ans = 0;
    for (int i = LG - 1; i >= 0; --i) {
      ans = max(ans, ans ^ basis[i]);
    }
    return ans;
  }
  bool can(T x) {
    for (int i = LG - 1; i >= 0; --i) {
      if (x >> i & 1) {
        if (!basis[i]) return false;
        x ^= basis[i];
      }
    }
    return true;
  }
  T kth(T k) {
    if (k < 1 || k > ((T)1 << sz)) return -1;
    T mask = 0;
    int64_t tot = 1LL << sz;
    for (int i = LG - 1; i >= 0; --i) {
      if (basis[i]) {
        int64_t low = tot / 2;
        if ((low < k && (~mask >> i & 1)) ||
        (low >= k && (mask >> i & 1))) {
          mask ^= basis[i];
        }
        if (low < k) k -= low;
        tot /= 2;
      }
    }
    return mask;
  }
\end{lstlisting}

\begin{lstlisting}
  // number of subsets having xor < x
  T count_less(T x) {
    if (x < 0) {
      return 0;
    }
    T ret = 0;
    T cnt = ((T)1 << sz);
    T mask = 0;
    for (int i = LG - 1; i >= 0; i--) {
      if (basis[i]) {
        if (x >> i & 1) {
          ret += (cnt >> 1);
          if (~mask >> i & 1) {
            mask ^= basis[i];
          }
        } else {
          if (mask >> i & 1) {
            mask ^= basis[i];
          }
        }
        cnt >>= 1;
      } else {
        if ((x ^ mask) >> i & 1) {
          if (x >> i & 1) {
            return ret + cnt;
          } else {
            return ret;
          }
        }
      }
    }
    return ret;
  }
  T count_leq(T x) {
    return count_less(x + 1);
  }
  T jumpK(T val, T k) {
    T p = count_leq(val);
    return kth(p + k);
  }
};
\end{lstlisting}

\textbf{Sum Of Digits upto N}
\begin{lstlisting}
uint64_t sumDigitsUpTo(uint64_t n){
  if(n <= 0) return 0;
  uint64_t sum = 0;
  for(uint64_t f = 1; f <= n; f *= 10){
    uint64_t hi = n / (f * 10);
    uint64_t cur = (n / f) % 10;
    uint64_t lo = n % f;
    sum += hi * 45LL * f;
    sum += cur * (lo + 1);
    sum += f * (cur * (cur - 1) / 2);
  }
  return sum;
}
\end{lstlisting}

\textbf{Number of digits upto N}
\begin{lstlisting}
int64_t get(int64_t n) {
  int64_t base = 1, digits = 1, now = 0;
  while (base * 10 <= n) {
    base *= 10;
    now += (base - (base / 10 == 0 ? 1 : base / 10))
    * digits;
    digits += 1;
  }
  n -= (base - 1);
  now += n * digits;
  return now;
}
\end{lstlisting}

\textbf{Linear Combination Checking}
Given \texttt{arr[] = \{a, b, c, d, ...\}} and an integer N, the task is to check if N can be represented as a linear combination of the array elements.

For example, $N = a \cdot x + b \cdot y + c \cdot z + \ldots$ where $x, y, z, \ldots$ are integers.

Ans: if $N \bmod \text{GCD}(\text{all array elements})$ is zero.
